%% ma: L11-B17-0069
%% lop: 11
%% chuong: 5
%% bai: 17
%% dang: 11.17.4
%% loai: DS
%% muc: [TH, TH, TH, TH]
%% dap_an: "SSĐS"
%% nguon: "tai-lieu-hoc-tap-mon-toan-lop-11-theo-chuong-trinh-giao-duc-pho-thong-2018.pdf (Lớp Toán Thầy Duy, Phạm Lê Duy, THPT Chu Văn An), Chương 3 Bài 3, Phần 2, Câu 9 (tr.417)"
%% kien_thuc: "Lớp 11 Bài 17 (hàm số liên tục)"
%% trang_thai: chinh-thuc
%% ghi_chu: "Đã duyệt 03/11/2026."
\begin{ex}
Cho hàm số $f(x)=\begin{cases}\dfrac{x^2}x&\text{khi }x<1,\ x\ne0\\0&\text{khi }x=0\\\sqrt x&\text{khi }x\ge1\end{cases}$. Khi đó
\choiceTF
{Hàm số liên tục tại mọi điểm trừ các điểm thuộc đoạn $[0;1]$.}
{Hàm số liên tục tại mọi điểm trừ điểm $x=0$.}
{\True Hàm số liên tục tại mọi điểm thuộc $\mathbb R$.}
{Hàm số liên tục tại mọi điểm trừ điểm $x=1$.}
\loigiai{\sloppy Với $x<1$, $x\ne0$: $f(x)=x$. Tại $x=0$: $\lim\limits_{x\to0}x=0=f(0)$, liên tục. Tại $x=1$: $\lim\limits_{x\to1^-}x=1$, $\lim\limits_{x\to1^+}\sqrt x=1$, $f(1)=1$, liên tục. Với $x>1$ hàm $\sqrt x$ liên tục. Vậy $f$ liên tục trên $\mathbb R$: chỉ ý c) đúng.}
\end{ex}
